\chapter{Quant bancaire}\label{ch:in:bank-quant}

Le 4 avril 2011, la Réserve fédérale et l'Office of the Comptroller of the Currency ont indiqué aux banques américaines que «~validation
doit être menée par des personnes qui ne sont responsables ni du développement ni de l'utilisation, et qui n'ont pas d'enjeu dans le fait qu'un modèle soit
déclaré valide~». Le régulateur du Royaume-Uni a inscrit le même principe dans sa propre déclaration, en vigueur depuis
le 17 mai 2024, et le guide de la Banque centrale européenne sur les modèles internes précise à quelle distance les deux groupes doivent se tenir. Le 17
avril 2026, les agences américaines ont remplacé leurs recommandations de 2011 par un texte plus court, fondé sur les risques, qui conserve l'idée. Chaque modèle
qu'une banque utilise a désormais un propriétaire, un développeur et un validateur, et l'intitulé du poste d'un quant bancaire indique en général lequel des trois
rôles occupe le quant. Ce chapitre décrit les quatre métiers de quant bancaire et le quant chargé des ajustements de valorisation qui se situe entre eux,
vérifie une ligne de reporting au regard des tests des superviseurs, et estime combien de validateurs
l'inventaire des modèles d'une banque emploie.

\begin{center}\footnotesize
\begin{tabular}{@{}p{2.3cm}p{2.5cm}p{2.5cm}p{2.5cm}p{2.5cm}@{}}
\toprule
\multicolumn{5}{@{}l}{\textbf{Fiches de rôle~: les quatre quants bancaires}}\\
\midrule
 & \omterm{def:in:bank-quant:strat}{stratégiste de desk} & \omterm{def:in:bank-quant:library}{quant de bibliothèque} & \omterm{def:in:bank-quant:risk}{quant de risque} & \omterm{def:in:bank-quant:validator}{validateur de modèles}\\
\midrule
est rattaché à & un desk de trading & un groupe central d'analyse & la fonction de risque & la gestion du risque de modèle\\
fournit & outils de valorisation et de risque pour le desk & la bibliothèque de valorisation & modèles de risque et de fonds propres & rapports de validation et constats\\
P\&L & soutient celui du desk & aucun & aucun & aucun\\
rend compte à & responsable du desk & responsable de l'analyse & directeur des risques (CRO) & responsable du risque de modèle\\
codes dans les déclarations & 13-2099.01, 15-2041 & 13-2099.01, 15-1252 & 13-2054, 13-2099.01 & 13-2054, 13-2099.01, 15-2041\\
traité dans & livre~5, chap.~27~; livre~9, chap.~24, 28 & livre~5, chap.~28~; livre~6, chap.~29 & livre~6, chap.~21, 23 & livre~6, chap.~26~; livre~12, chap.~21\\
\bottomrule
\end{tabular}
\end{center}

\section{Stratégiste de desk}

\begin{definition}[Stratégiste de desk]\label{def:in:bank-quant:strat}
Un \emph{stratégiste de desk}\index{stratégiste de desk} est un analyste quantitatif rattaché à l'un des desks de trading d'une banque et
qui construit les modèles, les outils d'analyse et les outils que les traders et les vendeurs du desk utilisent pour valoriser, couvrir et gérer leurs
positions et répondre aux clients.
\end{definition}

Les banques ne donnent pas le même nom à ce métier. Goldman Sachs appelle son personnel quantitatif des strats et les décrit dans une offre d'emploi~:
«~les stratégistes quantitatifs sont à la pointe de nos activités, et résolvent des problèmes concrets par une variété de
méthodes analytiques~»~; «~les strats de desk siègent sur les parquets de trading et apportent de la valeur aux clients internes \& externes grâce à
des modèles de pointe, des analyses prédictives et des outils de trading de haute qualité.~» D'autres banques parlent de quant de desk, de quant de front office
ou d'analyste quantitatif sur le desk~; le travail est le même.

Le contenu du travail dépend du desk.
\begin{itemize}
\item \textbf{Sur un desk d'options ou de produits structurés} (livre~5, chapitre~27~; livre~9, chapitre~28) le stratégiste valorise
de nouveaux payoffs, explique aux traders le risque du desk et repère les endroits où le modèle du desk ne capte pas un risque que porte le portefeuille.
\item \textbf{Sur un desk de flux} (livre~9, chapitre~24) le stratégiste construit les outils de cotation et de couverture du market
making électronique~: des pricers qui répondent aux demandes des clients en quelques millisecondes, des ajustements de cotation selon l'inventaire, des règles de couverture.
\item \textbf{Sur un central risk book} (livre~9, chapitre~25) le stratégiste compense les flux de plusieurs desks et décide
de ce qu'il faut internaliser et de ce qu'il faut couvrir à l'extérieur.
\end{itemize}
Un \omterm{def:in:bank-quant:strat}{stratégiste de desk} ne possède pas le résultat du desk, mais travaille au plus près de celui-ci~: une transaction mal valorisée ou une couverture manquante
apparaît dans l'explication du résultat du lendemain (livre~6, chapitre~27). Les retours sur ce poste sont rapides, et sa rémunération tend
à suivre les résultats annuels du desk (chapitre~13).

\section{Quant de bibliothèque}

\begin{definition}[Quant de bibliothèque]\label{def:in:bank-quant:library}
Un \emph{quant de bibliothèque}\index{quant de bibliothèque} est un analyste quantitatif qui construit et maintient la bibliothèque commune de valorisation et
de risque~: les méthodes numériques, les modèles et les interfaces que les outils de front office, les systèmes de risque et les contrôles de valorisation
appellent, dans des versions gérées sous contrôle de version et testées contre des indices de référence indépendants.
\end{definition}

Une banque qui valorise le même swap dans son système de trading, dans son moteur de risque et dans sa valorisation de fin de mois veut une seule
implémentation du swap, et non trois qui divergent. La bibliothèque est cette implémentation~: la construction des courbes, les modèles
d'options, les moteurs de Monte Carlo et d'arbres, les routines de calibration et les structures de données qui décrivent les transactions (livre~5,
chapitre~28, en construit une petite~; livre~6, chapitre~29, fait tourner un moteur de risque dessus). QuantLib, «~une bibliothèque libre et open source
pour la finance quantitative~», montre la forme d'une telle bibliothèque en accès public~; celle d'une banque, privée, est bien plus grande.

Le produit du \omterm{def:in:bank-quant:library}{quant de bibliothèque} est un logiciel, et ce rôle se situe entre le \omterm{def:in:bank-quant:strat}{stratégiste de desk} et le développeur quantitatif du
chapitre~19~: le quant choisit les mathématiques et écrit une bonne part du code numérique~; les développeurs construisent les systèmes autour
d'elle. Ce que livre le \omterm{def:in:bank-quant:library}{quant de bibliothèque} atteint tous les desks à la fois, si bien que le travail repose sur des mises en production, des tests de régression et
des comparaisons aux indices de référence plutôt que sur la demande du trader du matin~; et chaque changement apporté à un modèle utilisé par la banque passe
aux validateurs avant sa mise en production.

\section{Quant de risque}

\begin{definition}[Quant de risque]\label{def:in:bank-quant:risk}
Un \emph{quant de risque}\index{quant de risque} est un analyste quantitatif de la fonction de risque d'une banque qui construit les modèles
mesurant le risque et le capital de la banque~: la valeur en risque (VaR) et l'expected shortfall (ES), les scénarios de stress, l'exposition de contrepartie,
les paramètres de risque de crédit et les modèles relevant de l'approche des modèles internes (IMA).
\end{definition}

Le \omterm{def:in:bank-quant:risk}{quant de risque} travaille pour le directeur des risques (CRO), et non pour un desk, et ses modèles servent aux limites, aux fonds propres et à la vision du
conseil d'administration sur le risque plutôt que sur les prix. Le travail est enseigné dans tout le livre~6~: les mesures du risque de marché (chapitre~21), les tests de résistance
(chapitre~22), le capital des portefeuilles de négociation sous l'approche des modèles internes (chapitre~23), l'exposition de contrepartie
(chapitre~17) et les modèles de marge (chapitre~25). Deux traits la distinguent de celui du \omterm{def:in:bank-quant:strat}{stratégiste de desk}. D'abord, beaucoup de ses
modèles exigent l'approbation d'un superviseur, et en modifier un est un événement réglementaire, non une mise en production. Ensuite, ses utilisateurs sont
les gestionnaires de risque, la finance et les superviseurs, si bien que le modèle doit être expliqué autant qu'il tourne. Une banque qui utilise l'approche des
modèles internes doit voir ses modèles de risque réussir les tests de backtesting et d'attribution du résultat du livre~6,
chapitre~23, desk par desk, et le \omterm{def:in:bank-quant:risk}{quant de risque} est la personne qui répond lorsqu'un desk échoue à l'un d'eux.

\section{Validateur de modèles}

\begin{definition}[Validateur de modèles]\label{def:in:bank-quant:validator}
Un \emph{validateur de modèles}\index{validateur de modèles} est un analyste quantitatif de la fonction de gestion du risque de modèle d'une banque qui
examine des modèles construits par d'autres, indépendamment de leurs développeurs et de leurs propriétaires~: solidité conceptuelle, implémentation,
données, performance et limites, et qui consigne les constats que les développeurs doivent résoudre avant ou après que le modèle est
approuvé pour usage.
\end{definition}

La validation des modèles, la gestion du risque de modèle, l'inventaire des modèles, la hiérarchisation des modèles (tiering) et la remise en cause effective (effective challenge) sont définis dans
le livre~6, chapitre~26, et le propriétaire du modèle dans le livre~12, chapitre~21. Cette section porte sur la personne. Le travail du validateur
consiste à redériver, réimplémenter et mettre à mal les modèles d'autres personnes~: construire un indice de référence indépendant d'un pricer et le comparer
sur une grille de paramètres d'entrée, à tester les résultats d'un modèle de risque face aux pertes réalisées, à lire la documentation et à repérer ce
qu'elle ne dit pas. Son produit est un rapport de validation, une évaluation de la hiérarchisation et une liste de constats assortis de délais. Certaines banques
recrutent les deux filières par un même programme jeunes diplômés~; l'une d'elles décrit ainsi ses programmes de finance quantitative~: «~On
contribuera à développer ou à valider des modèles mathématiques, des méthodologies et des outils utilisés dans toute l'entreprise.~»

Ce qui rend la position du validateur inhabituelle, c'est que les superviseurs en fixent eux-mêmes le lieu.
\begin{itemize}
\item \textbf{États-Unis, 2011--2026.} La recommandation de 2011 exigeait que les validateurs «~n'aient aucun intérêt à ce qu'un modèle soit
jugé valide~», ajoutait que «~si l'indépendance peut être étayée par la séparation des lignes hiérarchiques, elle doit être
appréciée à l'aune des actes et des résultats~», et demandait aux banques de revoir chaque modèle «~au moins une fois par an~».
\item \textbf{États-Unis, depuis le 17 avril 2026.} La recommandation révisée demande une remise en cause effective par des personnes
dotées d'une «~indépendance suffisante pour préserver leur objectivité, ainsi que de la position organisationnelle et de l'influence
nécessaires pour obtenir tout changement~», désigne comme conflit à gérer le «~désalignement des incitations entre lignes hiérarchiques
distinctes, telles que les groupes de développement et de validation des modèles~», et laisse le calendrier et la fréquence de la validation varier selon le modèle.
\item \textbf{Royaume-Uni, depuis le 17 mai 2024.} Les établissements disposant d'une approbation pour modèles internes doivent
«~démontrer leur indépendance par des lignes hiérarchiques distinctes pour les \omterm{def:in:bank-quant:validator}{validateurs de modèles} et pour les développeurs et propriétaires de modèles~».
\item \textbf{Zone euro.} Le guide de la BCE énumère trois dispositifs, du plus robuste au moins robuste, et attend des établissements
grands et complexes qu'ils recourent au premier (la \cref{fig:in:bank-quant:org}).
\end{itemize}

\begin{definition}[Dispositif d'indépendance]\label{def:in:bank-quant:arrangement}
Pour un validateur et un développeur, les dispositifs du guide de la BCE sont~: (a) deux unités rattachées à des membres
différents de la haute direction~; (b) deux unités rattachées au même membre de la haute direction~; (c) des membres du personnel distincts au sein d'une même unité.
Un validateur qui rend compte, directement ou par l'intermédiaire d'un manager, au développeur ou au propriétaire du modèle n'est pas du tout indépendant.
\end{definition}

\begin{dated}{2026-09}{Les règles de risque de modèle qui encadrent le travail du validateur}\label{dat:in:bank-quant:rules}
États-Unis~: la SR~11-7 (4 avril 2011) et le OCC Bulletin 2011-12 ont été remplacés le 17 avril 2026 par la recommandation révisée
de la Réserve fédérale, de l'OCC et de la FDIC sur la gestion du risque de modèle (SR~26-2~; OCC Bulletin 2026-13), «~censée
concerner au premier chef les organisations bancaires dont le total de bilan dépasse 30~milliards de dollars~»~; elle ne fixe aucune «~norme contraignante ni exigence prescriptive
»~et laisse en dehors de son champ les modèles d'IA générative et agentique. Royaume-Uni~: la PRA SS1/23, en vigueur depuis
le 17 mai 2024, pour les banques et les entreprises d'investissement disposant d'une approbation pour modèles internes~; hiérarchisation selon la matérialité et la complexité~; lignes hiérarchiques distinctes
pour les validateurs. Zone euro~: guide de la BCE sur les modèles internes, version 4.1 du 26 juin 2026 (version 4.0 du 28
juillet 2025)~: une validation interne initiale puis annuelle de chaque modèle interne~; dispositifs (a), (b) et (c).
\end{dated}

\begin{omfigure}
\begin{tikzpicture}[font=\scriptsize,
  box/.style={draw,rounded corners=2pt,align=center,minimum height=0.8cm,text width=1.75cm,inner sep=2pt},
  val/.style={box,fill=omDef!15},dev/.style={box,fill=omThm!12},
  pair/.style={dashed,thick,omThm},lab/.style={fill=white,inner sep=1pt,font=\scriptsize\bfseries}]
\node[box] (ceo) at (-0.5,0) {directeur\\général};
\node[box] (cro) at (-3.7,-1.3) {directeur des\\risques (CRO)};
\node[box] (mkt) at (2.6,-1.3) {responsable des\\marchés};
\node[box] (mrm) at (-5.25,-2.7) {responsable du\\risque de modèle};
\node[box] (ra) at (-2.1,-2.7) {responsable de\\l'analyse des risques};
\node[box] (qa) at (1.05,-2.7) {responsable de\\l'analyse quant};
\node[box] (desk) at (4.2,-2.7) {responsable du\\desk};
\node[val] (v1) at (-5.25,-4.2) {validateur};
\node[dev] (rd) at (-3.15,-4.2) {développeur\\risque};
\node[val] (rv) at (-1.05,-4.2) {validateur\\risque};
\node[dev] (ld) at (1.05,-4.2) {développeur\\de bibliothèque};
\node[dev] (ds) at (3.15,-4.2) {stratégiste\\de desk};
\node[val] (dv) at (5.25,-4.2) {validateur\\de desk};
\foreach \a/\b in {ceo/cro,ceo/mkt,cro/mrm,cro/ra,mkt/qa,mkt/desk,mrm/v1,ra/rd,ra/rv,qa/ld,desk/ds,desk/dv}
  \draw (\a.south) -- (\b.north);
\draw[pair] ([xshift=4mm]v1.south) to[out=-60,in=-120] node[lab,pos=0.5] {(b)} ([xshift=-4mm]rd.south);
\draw[pair] ([xshift=4mm]rd.south) to[out=-60,in=-120] node[lab,pos=0.5] {(c)} ([xshift=-4mm]rv.south);
\draw[pair] ([xshift=-2mm]v1.south) to[out=-70,in=-110] node[lab,pos=0.5] {(a)} (ld.south);
\draw[pair] ([xshift=4mm]ds.south) to[out=-60,in=-120] node[lab,pos=0.5] {échec} ([xshift=-4mm]dv.south);
\end{tikzpicture}

\omcaption{Une organisation bancaire illustrative et quatre couples validateur--développeur (en pointillés), classés par
\texttt{firm.roles.independence}~: (a) membres différents de la haute direction~; (b) unités différentes sous un même
responsable de haut niveau~; (c) membres du personnel distincts au sein d'une même unité~; «~échec~»~: le validateur rend compte au propriétaire du modèle qu'il examine.
La haute direction est ici représentée par le directeur des risques et le responsable des marchés. Données~: \texttt{in\_bankquant.org\_cases}.}\label{fig:in:bank-quant:org}
\end{omfigure}

\section{Quant de l'ajustement de valeur}

Le desk qui gère les ajustements de valeur de crédit, de financement et de fonds propres d'une banque (livre~6, chapitre~20)
a besoin de quants qui lui sont propres, et ceux-ci ne rentrent pas proprement dans les quatre fiches. Ils construisent des simulations d'exposition sur l'ensemble du portefeuille, comme un \omterm{def:in:bank-quant:risk}{quant de risque}~;
ils valorisent les ajustements facturés aux desks de trading, comme un \omterm{def:in:bank-quant:strat}{stratégiste de desk}~; leurs chiffres entrent dans les comptes de la banque, si bien que le product control (livre~6, chapitre~27)
et les validateurs les examinent. Le poste exige la modélisation la plus large parmi les quants de banque
(la dynamique de chaque classe d'actifs dans une même simulation) et le plus de calcul, et c'est là qu'un \omterm{def:in:bank-quant:risk}{quant de risque} passe le plus souvent
au front office, ou l'inverse (chapitre~28).

\section{À qui ils rendent compte et ce qu'ils livrent}

Les quatre fiches diffèrent surtout par la personne à qui le quant rend compte, et c'est elle qui décide de ce qui compte comme un bon travail.
\begin{center}\footnotesize
\begin{tabular}{@{}p{2.6cm}p{4.9cm}p{5.1cm}@{}}
\toprule
rôle & jugé sur & ce qui tourne mal\\
\midrule
\omterm{def:in:bank-quant:strat}{stratégiste de desk} & l'usage des outils par le desk~; valorisation et couverture qui tiennent la route & un modèle qui convient mieux au desk qu'au risque\\
\omterm{def:in:bank-quant:library}{quant de bibliothèque} & exactitude, rapidité, stabilité d'une version à l'autre & un changement qui modifie d'un coup les chiffres de tous les desks\\
\omterm{def:in:bank-quant:risk}{quant de risque} & fonds propres, limites et tests prudentiels réussis & un modèle calibré pour le test plutôt que pour le risque\\
\omterm{def:in:bank-quant:validator}{validateur de modèles} & constats pertinents, remontés à temps & une revue qui valide ce qu'on lui a montré\\
quant XVA & ajustements conformes à ceux du marché et à ceux des comptes & un modèle d'exposition trop lent pour tourner chaque jour\\
\bottomrule
\end{tabular}
\end{center}

Le nombre de validateurs est un choix de conception pour la banque, et un choix lourd. Il découle de l'inventaire des modèles, de leur
hiérarchisation et du cycle de revue que la politique de la banque fixe pour chaque niveau.

\begin{method}[Charge de validation]\label{met:in:bank-quant:workload}
Soit $n_t$ le nombre de modèles de niveau $t$ dans l'inventaire~; une validation complète d'un modèle de niveau $t$ prend $F_t$ heures et une revue périodique
$R_t$ heures~; le niveau $t$ fait l'objet d'une revalidation complète tous les $k_t$ ans et d'une revue les autres années~; et une part $c$ de l'ensemble des modèles
change de manière significative chaque année et fait l'objet d'une nouvelle validation complète. Les heures annuelles sont
\[
H=\sum_t n_t\left(\frac{F_t}{k_t}+R_t\Bigl(1-\frac{1}{k_t}\Bigr)+c\,F_t\right),
\]
et le nombre de validateurs nécessaires est $H/h$, pour $h$ heures productives par validateur et par an.
\end{method}

\begin{example}[Un inventaire de 1\,000 modèles]\label{ex:in:bank-quant:inventory}
Attribuez à chaque modèle une note de 1 à 3 pour la matérialité, la complexité et l'incertitude, chaque note étant tirée avec les probabilités 0.5, 0.3 et
0.2, puis hiérarchisez-le avec la règle du livre~6 (deux fois la note de matérialité plus les deux autres~: niveau~1 à partir de 10, niveau~2 à partir de 7, niveau~3 en dessous), ce qui fixe
une revalidation tous les ans, tous les deux ans ou tous les trois ans, selon le niveau. L'inventaire attendu compte 102 modèles de niveau~1, 403 de niveau~2 et 495 de niveau~3.
Avec les heures illustratives $F=(400,160,60)$, $R=(40,16,8)$, $c=0.10$ et $h=1\,600$, l'inventaire demande
102\,302 heures par an, soit 63.9 validateurs. Les modèles de niveau~1 représentent 10.2~\% de l'inventaire et absorbent 43.9~\% des heures.
\end{example}

\begin{center}\footnotesize
\begin{tabular}{@{}lrrr@{}}
\toprule
niveau & modèles & heures par an & part des heures\\
\midrule
1 & 102 & 44\,880 & 43.9~\%\\
2 & 403 & 41\,912 & 41.0~\%\\
3 & 495 & 15\,510 & 15.2~\%\\
\bottomrule
\end{tabular}
\end{center}

\begin{omfigure}
\begin{tikzpicture}
\begin{axis}[width=10.5cm,height=5.4cm,xmin=0,xmax=45,ymin=40,ymax=150,
  xlabel={\small part des modèles de niveau~1 sur 1\,000,~\%},ylabel={\small validateurs nécessaires},
  tick label style={font=\footnotesize},grid=major,grid style={gray!20},table/col sep=comma,
  legend style={font=\scriptsize,at={(0.03,0.97)},anchor=north west,legend cell align=left}]
\addplot[omDef,thick,mark=*] table[x=tier1_pct,y=validators]{figdata/industry/18-bank-quant/headcount.csv};
\addplot[only marks,mark=square*,mark size=2.5pt,omThm] table[x=tier1_pct,y=validators]{figdata/industry/18-bank-quant/base.csv};
\legend{effectif de niveau~2 fixé à 403,inventaire de l'exemple}
\end{axis}
\end{tikzpicture}

\omcaption{Validateurs nécessaires pour un inventaire de 1\,000 modèles à mesure que la part de niveau~1 augmente (les modèles de niveau~3 deviennent
niveau~1), avec les heures et les cycles de l'exemple. Les heures sont illustratives~; la pente, d'environ 2.6 validateurs par
point de pourcentage de modèles de niveau~1, est le coût de la décision de hiérarchisation. Données~: \texttt{in\_bankquant.curve}.}\label{fig:in:bank-quant:headcount}
\end{omfigure}

Deux conclusions résistent à tout choix raisonnable des heures. La règle de hiérarchisation fixe la taille de la fonction de validation~: déplacer
un modèle du niveau~3 au niveau~1 ajoute 409 heures par an, soit un quart de validateur. Et la recommandation américaine révisée, qui
laisse la fréquence de validation varier selon le modèle, déplace l'effort des nombreux petits modèles vers les quelques grands~;
la règle de la zone euro, qui impose une validation annuelle de chaque modèle interne, fait l'inverse pour les modèles de capital.

\section{Rémunération~: ce que montrent les dossiers des quants de banque}

Les labour condition applications du chapitre~14 ne répartissent pas les quants de banque dans les quatre familles~: les intitulés de poste varient trop
d'une banque à l'autre, et la famille «~\omterm{def:in:quant-researcher:def}{chercheur quantitatif} ou analyste~» des banques en regroupe la plupart, avec une famille distincte
pour les intitulés de risque. L'enquête sur les professions ajoute les spécialistes du risque financier (13-2054) par secteur.

\begin{dated}{2025-09}{Ce que gagnent les quants de banque~: les données publiques}\label{dat:in:bank-quant:pay}
Labour condition applications américaines des neuf employeurs bancaires identifiés, \omterm{def:in:how-pay-works:base}{salaire de base} proposé~: famille «~\omterm{def:in:quant-researcher:def}{chercheur quantitatif} ou
analyste~», médiane de 128\,300~\$ pour l'exercice 2021 (895 dossiers) et 158\,100~\$ pour l'exercice 2025 (825)~; par \omterm{def:in:pay-levels-by-role-firm-type-and-seniority:soc}{niveau de salaire}
pour l'exercice 2025~: 88\,300~\$ (I), 145\,300~\$ (II), 179\,335~\$ (III), 200\,000~\$ (IV). Intitulés de risque, exercice 2025~: médiane
136\,776~\$ (380 dossiers, 5 employeurs). Enquête sur les professions, mai 2025, spécialistes du risque financier~: banques (intermédiation
de crédit), médiane 108\,170~\$, 10\up{e}--90\up{e}~centile 62\,620--187\,490~\$ (17\,810 salariés)~; secteur des valeurs mobilières, médiane
133\,070~\$, 83\,230--219\,990~\$ (9\,840). \omterm{def:in:how-pay-works:base}{Salaire de base} uniquement.
\end{dated}

\begin{omfigure}
\begin{tikzpicture}
\begin{axis}[width=9.2cm,height=5.4cm,xmin=40,xmax=240,ymin=0.4,ymax=5.6,ytick=data,
  yticklabels={{enquête~: spécialistes du risque~; banques},{enquête~: spécialistes du risque~; valeurs mobilières},{enquête~: autres spécialistes financiers~; valeurs mobilières},{dépôts~: postes de risque en banque},{dépôts~: chercheurs quant en banque}},yticklabel style={font=\scriptsize},
  xticklabel style={font=\footnotesize},xlabel={\small salaire de base annuel, milliers de dollars},grid=major,grid style={gray!20},
  table/col sep=comma]
\addplot[draw=none,mark=none,error bars/.cd,x dir=both,x explicit,error mark=none,error bar style={omDef,line width=1pt}]
  table[x=p50,y=pos,x error plus expr=\thisrow{p90}-\thisrow{p50},x error minus expr=\thisrow{p50}-\thisrow{p10}]
  {figdata/industry/18-bank-quant/ranges.csv};
\addplot[draw=none,mark=none,error bars/.cd,x dir=both,x explicit,error mark=none,error bar style={omDef!40,line width=7pt}]
  table[x=p50,y=pos,x error plus expr=\thisrow{p75}-\thisrow{p50},x error minus expr=\thisrow{p50}-\thisrow{p25}]
  {figdata/industry/18-bank-quant/ranges.csv};
\addplot[only marks,mark=|,mark size=5pt,line width=1.5pt,omThm] table[x=p50,y=pos]{figdata/industry/18-bank-quant/ranges.csv};
\end{axis}
\end{tikzpicture}

\omcaption{\omterm{def:in:how-pay-works:base}{Salaire de base} des quants de banque, selon deux sources~: le \omterm{def:in:how-pay-works:base}{salaire de base} proposé dans les labour condition applications de l'exercice 2025 («~dossiers~»)
et l'enquête sur les professions de mai 2025 («~enquête~»), du 10\up{e} au 90\up{e}~centile (trait fin), écart interquartile (trait épais), médiane
(repère). La catégorie 13-2099 de l'enquête est plus large que les analystes quantitatifs qu'elle contient. Données~:
\texttt{data/industry/lca\_ranges.csv} et \texttt{oews\_roles.csv}, via \texttt{in\_bankquant}.}\label{fig:in:bank-quant:pay}
\end{omfigure}

Les dossiers se situent au-dessus de l'enquête pour la même raison qu'au chapitre~14~: ils couvrent les employeurs qui déposent des demandes
d'embauche de salariés étrangers qualifiés, au salaire qu'elles proposent, alors que l'enquête couvre l'ensemble de la profession. Entre l'exercice 2021 et 2025, la
médiane de l'offre des banques pour la famille quantitative a progressé de 23.2~\% en dollars et de 3.7~\% après les prix à la consommation américains, et l'échelle
par niveau s'est accentuée (la \cref{fig:in:bank-quant:levels})~: les niveaux II à IV ont augmenté, le niveau~I non.

\begin{omfigure}
\begin{tikzpicture}
\begin{axis}[width=10.5cm,height=5cm,xmin=0.8,xmax=4.2,ymin=60,ymax=220,xtick={1,2,3,4},xticklabels={I,II,III,IV},
  xlabel={\small niveau de salaire},ylabel={\small base médiane, milliers de \$},tick label style={font=\footnotesize},
  grid=major,grid style={gray!20},table/col sep=comma,unbounded coords=jump,
  legend style={legend cell align=left,font=\scriptsize,at={(0.5,-0.3)},anchor=north,/tikz/every even column/.append style={column sep=8pt}},
  legend columns=3]
\addplot[black!60,thick,dashed,mark=o,mark options={solid}] table[x=level,y=qr2021]{figdata/industry/18-bank-quant/levels.csv};
\addplot[omDef,thick,mark=*] table[x=level,y=qr2025]{figdata/industry/18-bank-quant/levels.csv};
\addplot[omThm,thick,mark=square*] table[x=level,y=risk2025]{figdata/industry/18-bank-quant/levels.csv};
\legend{famille quantitative 2021,famille quantitative 2025,intitulés de risque 2025}
\end{axis}
\end{tikzpicture}

\omcaption{\omterm{def:in:how-pay-works:base}{Salaire de base} proposé médian dans les banques par \omterm{def:in:pay-levels-by-role-firm-type-and-seniority:soc}{niveau de salaire}~: la famille quantitative aux exercices 2021 et 2025, et les intitulés de risque
à l'exercice 2025 (les cellules de risque 2021 sont trop petites au niveau~I). Données~: comme pour la \cref{fig:in:bank-quant:pay}.}\label{fig:in:bank-quant:levels}
\end{omfigure}

\section{Tutoriel~: rémunération, lignes hiérarchiques et charge de travail des validateurs}\label{tut:in:bank-quant:tut}

\begin{tutorial}
\textbf{Objectif.} Placer les éléments de rémunération du quant de banque à côté d'une vérification des lignes hiérarchiques et de la taille de la
fonction de validation. \textbf{État final~:} les \cref{fig:in:bank-quant:org,fig:in:bank-quant:headcount,fig:in:bank-quant:pay,fig:in:bank-quant:levels}.
\begin{tutsteps}
\item \textbf{Rémunération.} \texttt{firm.roles.lca\_cells(rows, role, fy, level)} pour les familles \texttt{quant researcher} et
\texttt{risk}, en gardant les cellules bancaires~; \texttt{data/industry/oews\_roles.csv} pour la profession 13-2054 dans les secteurs
5220A1 et 523000 (dérivé une fois par \texttt{in\_oews\_roles\_derive.py} à partir des fichiers de l'enquête).
\item \textbf{Lignes hiérarchiques.} Écrire l'organisation comme une carte reliant chaque personne ou unité à son responsable, et marquer les
membres de la direction générale~; \texttt{independence} attribue une note à chaque couple validateur--développeur
(le \cref{lst:in:bank-quant:independence}).
\omcode{code/firm/roles/firm_roles.py}{165}{183}{Indépendance d'un validateur vis-à-vis du développeur et du propriétaire d'un modèle, notée selon
les dispositions du guide de la BCE.}\label{lst:in:bank-quant:independence}
\item \textbf{Inventaire.} \texttt{firm.modelval.ModelRecord} du livre~6 hiérarchise chaque modèle à partir de ses notes et fixe son
intervalle de revalidation~; \texttt{in\_bankquant.inventory} prend l'effectif attendu de chaque niveau.
\item \textbf{Charge de travail.} \texttt{validation\_hours} et \texttt{validator\_headcount} appliquent la méthode
(le \cref{lst:in:bank-quant:hours})~; \texttt{in\_bankquant.curve} fait varier la part de niveau~1.
\omcode{code/firm/roles/firm_roles.py}{186}{198}{Heures de validation par an et validateurs nécessaires.}\label{lst:in:bank-quant:hours}
\end{tutsteps}
\end{tutorial}

Pour l'exemple~: 63.9 validateurs~; 50.7 avec une part de niveau~1 de 5~\% et 140.1 avec 40~\%.

\textbf{Pour aller plus loin.} Attribuer à chaque niveau sa propre part de changements substantiels~; ajouter la validation des modèles d'apprentissage automatique avec
la liste de contrôle du livre~12 (\texttt{firm.modelcard.validation\_checklist}), dont les tests d'explication et de stabilité prennent plus de temps que le benchmark d'un pricer~;
tirer l'inventaire au hasard plutôt que de prendre son espérance.

\section{Construction~: fiches de quant bancaire, lignes hiérarchiques et charge de validation}\label{bld:in:bank-quant:roles}

\begin{build}
\bfield{Rôle} Ajouter les quatre quants bancaires au registre des rôles, une vérification des lignes hiérarchiques et le modèle de la
taille de la fonction de validation.
\bfield{Interface} \texttt{firm.roles}~: les fiches \texttt{desk strategist}, \texttt{library quant}, \texttt{risk quant},
\texttt{model validator}~; \texttt{chain(boss, node)}~; \texttt{independence(boss, senior, validator, developer, owner)}
renvoie \texttt{a}, \texttt{b}, \texttt{c} ou \texttt{fail}~;
\texttt{validation\_hours(counts, full\_hours, review\_hours, interval\_years, change\_rate)}~; \texttt{validator\_headcount(\dots, productive\_hours)}.
\bfield{Règles} Un cycle de rattachement est une erreur~; un validateur placé sous le développeur ou le propriétaire échoue quelles que soient les unités~; les heures
et les cycles sont ceux de l'appelant, signalés comme illustratifs~; les niveaux proviennent de la règle du livre~6.
\bfield{Tests d'acceptation} \texttt{code/firm/roles/tests/}~: les quatre grades sur une petite organisation~; un cycle lève une exception~; un
niveau avec un cycle d'un an coûte exactement sa validation complète~; un cycle de quatre ans coûte un quart de celle-ci plus trois quarts
d'une revue.
\bfield{Pour aller plus loin} Propriétaires et développeurs comme ensembles (un modèle avec plusieurs développeurs)~; une organisation matricielle avec deux managers
par personne~; les constats et leurs échéances comme une file d'attente dont l'effectif doit résorber l'arriéré.
\end{build}

\begin{omsources}
\item Board of Governors of the Federal Reserve System et OCC (2011), SR~11-7, Supervisory guidance on model risk
management.
\item Federal Reserve, OCC et FDIC (2026), SR~26-2, Revised guidance on model risk management~; OCC Bulletin 2026-13.
\item Prudential Regulation Authority (2023), SS1/23, Model risk management principles for banks.
\item Banque centrale européenne (2026), ECB guide to internal models, release 4.1.
\item Pages carrières de Goldman Sachs et JPMorganChase~; QuantLib.
\item Tableaux du chapitre~14 issus des fichiers LCA du Department of Labor~; enquête sur les professions du BLS, mai 2025.
\end{omsources}

\section{Exercices}

\begin{exercise}[$\star$]\label{exo:in:bank-quant:1}
De combien la médiane des offres des banques pour la famille quantitative a-t-elle augmenté entre les exercices 2021 et 2025, en dollars et
après correction par les prix à la consommation américains (indice 114.325 en 2021, 135.8312 en 2025)~?
\end{exercise}

\begin{exercise}[$\star$]\label{exo:in:bank-quant:2}
Un validateur rend compte au directeur des risques par l'intermédiaire du responsable du risque de modèle~; le développeur du modèle rend compte au responsable
des marchés par l'intermédiaire du responsable de l'analyse quantitative. Quelle organisation est-ce~?
\end{exercise}

\begin{exercise}[$\star$]\label{exo:in:bank-quant:3}
Avec les paramètres de l'exemple, combien d'heures par an coûte un modèle de niveau~1, et un modèle de niveau~3~?
\end{exercise}

\begin{exercise}[$\star\star$]\label{exo:in:bank-quant:4}
Supposons que chaque modèle de niveau~2 soit un modèle de fonds propres qui doit être validé intégralement chaque année. Combien de validateurs
l'inventaire de l'exemple exige-t-il~?
\end{exercise}

\begin{exercise}[$\star\star$]\label{exo:in:bank-quant:5}
Un nouveau programme de produits double la part de modèles modifiés substantiellement chaque année, pour la porter à 20~\%. Combien de validateurs
faut-il~?
\end{exercise}

\begin{exercise}[$\star\star$]\label{exo:in:bank-quant:6}
Pourquoi le \omterm{def:in:bank-quant:strat}{stratégiste de desk} qui a construit un pricer ne doit-il pas le valider, même s'il est la personne la mieux qualifiée
de la banque~?
\end{exercise}

\begin{exercise}[$\star\star\star$]\label{exo:in:bank-quant:7}
\emph{Programmation.} Tirer au hasard les scores des 1\,000 modèles 2\,000 fois (graine 18) au lieu de prendre les effectifs espérés par niveau.
Quelle est la moyenne et l'intervalle entre le 5\up{e} et le 95\up{e}~centile des validateurs nécessaires~?
\end{exercise}

\begin{exercise}[$\star\star\star$]\label{exo:in:bank-quant:8}
\emph{Trouver la faille.} «~Nos validateurs forment leur propre équipe sous le responsable de l'analyse quantitative, dont ils
relisent la bibliothèque, et celui-ci rend compte au responsable des marchés~; c'est une unité séparée, donc il s'agit de l'organisation (b).~»
\end{exercise}

\section{Problème~: Qui valide les validateurs~?}

\begin{problem}[{Problème du week-end --- qui valide les validateurs~?}]\label{pb:in:bank-quant:1}
Le nouveau responsable du risque de modèle d'une banque doit proposer un budget pour la validation et défendre la ligne hiérarchique de son équipe.

\medskip\noindent\textbf{Partie I --- Les rôles.}
\begin{enumerate}
\item Définir le \omterm{def:in:bank-quant:strat}{stratégiste de desk}, le \omterm{def:in:bank-quant:library}{quant de bibliothèque}, le \omterm{def:in:bank-quant:risk}{quant de risque} et le \omterm{def:in:bank-quant:validator}{validateur de modèles}.
\item Comment une banque décrit-elle ses desk strats~?
\item Pourquoi une banque veut-elle une seule bibliothèque de pricing plutôt qu'une par système~?
\item Qu'est-ce qui fait d'une modification de modèle par un \omterm{def:in:bank-quant:risk}{quant de risque} un événement réglementaire~?
\item Où se situe le quant d'ajustements de valorisation, et pourquoi le poste ne correspond-il à aucune fiche unique~?
\end{enumerate}

\medskip\noindent\textbf{Partie II --- Les règles.}
\begin{enumerate}[resume]
\item Que disaient les orientations américaines de 2011 sur qui doit valider, et sur les lignes hiérarchiques~?
\item Qu'est-ce qui les a remplacées en 2026, et qu'est-ce qui a changé quant à la fréquence de validation~?
\item Qu'attend la déclaration britannique des entreprises disposant de l'agrément des modèles internes~?
\item Énoncer les trois dispositions de la BCE.
\item Selon les orientations américaines révisées, qui contrôle la fonction de risque de modèle elle-même, et que ne fait-il pas~?
\end{enumerate}

\medskip\noindent\textbf{Partie III --- La charge de travail.}
\begin{enumerate}[resume]
\item Énoncer la méthode de charge de validation.
\item Comment la règle de hiérarchisation du livre~6 transforme-t-elle les scores en niveaux et en cycles~?
\item Donner les effectifs espérés par niveau pour 1\,000 modèles avec les probabilités de score de l'exemple.
\item Quelle part des heures reviennent aux modèles de niveau~1, et pourquoi~?
\item Combien de validateurs une validation annuelle complète des modèles de niveau~2 ajoute-t-elle~?
\end{enumerate}

\medskip\noindent\textbf{Partie IV --- Le verdict.}
\begin{enumerate}[resume]
\item Énoncer le \emph{résultat nommé}~: les validateurs nécessaires pour 1\,000 modèles avec la répartition par niveau de l'exemple, et la variation
pour chaque point de pourcentage de part de niveau~1.
\item Lesquels des paramètres de l'exemple sont illustratifs, et lesquels proviennent de règles~?
\item Évaluer sa proposition de ligne~: son équipe rend compte au directeur des risques, les développeurs au responsable des marchés.
\item Qui valide les validateurs~?
\item En deux phrases, présenter son dossier budgétaire.
\end{enumerate}
\end{problem}

\section{Questions d'entretien}

\begin{interviewq}[$\star$ \iqroles{banque, risque}]\label{iq:in:bank-quant:1}
Comment testeriez-vous un nouveau pricer Black--Scholes que le desk a écrit~?
\end{interviewq}

\begin{interviewq}[$\star$ \iqroles{banque, trader}]\label{iq:in:bank-quant:2}
Un trader affirme que le prix de votre pricer pour une option à barrière est faux, parce qu'un concurrent cote différemment. Que faites-vous~?
\end{interviewq}

\begin{interviewq}[$\star\star$ \iqroles{risque}]\label{iq:in:bank-quant:3}
La valeur en risque (VaR) de votre desk a connu 7 dépassements en 250 jours à 99~\%. Le modèle est-il faux~?
\end{interviewq}

\begin{interviewq}[$\star\star$ \iqroles{banque, développeur}]\label{iq:in:bank-quant:4}
Une nouvelle version de la bibliothèque modifie une interpolation de courbe et déplace de quelques points de base la valorisation de chaque desk. Comment
la mettez-vous en production~?
\end{interviewq}

\begin{interviewq}[$\star\star$ \iqroles{banque, risque}]\label{iq:in:bank-quant:5}
Que regarderiez-vous en premier dans la documentation d'un modèle que vous n'avez jamais vu~?
\end{interviewq}

\begin{interviewq}[$\star\star\star$ \iqroles{banque, chercheur}]\label{iq:in:bank-quant:6}
Deux calibrations de volatilité locale ajustent aussi bien la surface d'aujourd'hui et évaluent un cliquet différemment de 3~\%. Laquelle
utilisez-vous, et que dites-vous aux validateurs~?
\end{interviewq}
